%==============================================================================
% RELATED WORK UPDATE - Adding PDB Validation Context
%==============================================================================

\section{Related Work}

\subsection{Glycan Structure Validation Methods}

% NEW: PDB-based validation context
Several approaches exist for validating glycan stereochemistry in predicted structures. 

\textbf{Privateer} \citep{agirre2015} is the standard tool for validating glycan 
structures in experimental PDB deposits, checking anomeric configuration, 
absolute configuration (D/L), ring puckering, and linkage geometry. However, 
Privateer requires 3D coordinates, making it unsuitable for validating input 
format specifications before structure prediction.

\textbf{PDB Structure Comparison} provides an alternative validation approach. 
By comparing tool output against experimentally verified PDB structures, 
stereochemistry correctness can be validated without running predictions. 
We employed this approach in our validation:

\begin{itemize}
    \item \textbf{PDB:2VXR} (Sialyllactose): Validates C2 anomeric position 
    for Neu5Ac---a critical test case where SMILES format fails
    \item \textbf{PDB:5NSC} (Fucosylated): Validates L-configuration and 
    alpha-anomer for fucose, a common error case from literature
    \item \textbf{PDB:5K65} (M3 N-glycan): Validates branch topology and 
    N-glycosidic linkage specification
    \item \textbf{PDB:1L6R} (Lactose): Baseline validation for epimer 
    distinction (galactose vs glucose)
\end{itemize}

\subsection{Glycan Databases and Standards}

The \textbf{Chemical Component Dictionary (CCD)} maintained by the Worldwide 
PDB provides standardized residue identifiers. Our CCD mapper uses CCD codes 
as the authoritative reference, ensuring compatibility with PDB infrastructure.

\textbf{GlyTouCan} \citep{tiemeyer2016} and \textbf{GlyGen} \citep{campbell2020} 
provide standardized glycan notations. GlycoSMILES2BAP integrates with these 
databases through IUPAC-condensed and WURCS format support.

\subsection{Comparison with Existing Tools}

\begin{table}[h]
\centering
\caption{Comparison of Glycan Validation Approaches}
\begin{tabular}{lcccc}
\toprule
Method & Input & Output & Validation Type & AF3 Ready \\
\midrule
Privateer & 3D coordinates & Validation report & Post-prediction & No \\
PDB Comparison & CCD codes & Accuracy metrics & Pre-prediction & Yes \\
GlycoSMILES2BAP & IUPAC/WURCS & CCD+BAP & Pre-prediction & Yes \\
\bottomrule
\end{tabular}
\end{table}

% NEW: Key differentiation
Our PDB structure comparison validation (Section~\ref{sec:pdb_validation}) 
demonstrates that GlycoSMILES2BAP correctly handles the critical stereochemistry 
cases identified in literature, particularly sialic acid C2 anomeric position 
and fucose L-configuration.
